\code{BPatch\_flowGraph} represents the control flow graph of a
function. It provides methods for discovering the
basic blocks and loops within the function (using which a caller can
navigate the graph). A \code{BPatch\_flowGraph} object
can be obtained by calling the getCFG method of a
\code{BPatch\_function} object.

\begin{apient}
  bool containsDynamicCallsites()
\end{apient}

\apidesc{
  Return true if the control flow graph contains any dynamic call
  sites (e.g., calls through a function pointer).
}

\begin{apient}
  void getAllBasicBlocks(std::set<BPatch_basicBlock*>&)
\end{apient}

\apidesc{
  Fill the given set with pointers to all basic blocks in the control
  flow graph.
}

\begin{apient}
  void getEntryBasicBlock(std::vector<BPatch_basicBlock*>&)
\end{apient}

\apidesc{
  Fill the given vector with pointers to all basic blocks that are
  entry points to the function.
}

\begin{apient}
  void getExitBasicBlock(std::vector<BPatch_basicBlock*>&)
\end{apient}

\apidesc{
  Fill the given vector with pointers to all basic blocks that are
  exit points of the function.
}

\begin{apient}
  void getLoops(std::vector<BPatch_basicBlockLoop*>&)
\end{apient}

\apidesc{
  Fill the given vector with a list of all natural (single entry)
  loops in the control flow graph.
}

\begin{apient}
  void getOuterLoops(std::vector<BPatch_basicBlockLoop*>&)
\end{apient}

\apidesc{
  Fill the given vector with a list of all natural (single entry)
  outer loops in the control flow graph.
}

\begin{apient}
  BPatch_loopTreeNode *getLoopTree()
\end{apient}

\apidesc{
  Return the root node of the tree of loops in this flow graph.
}

\begin{apient}
  enum BPatch_procedureLocation {
    BPatch_locLoopEntry,
    BPatch_locLoopExit,
    BPatch_locLoopStartIter,
    BPatch_locLoopEndIter
  }
\end{apient}

\apidesc{
  The place in the loop structure of the current procedure.
}

\begin{apient}
  std::vector<BPatch_point*> *findLoopInstPoints(const BPatch_procedureLocation loc,
                                                 BPatch_basicBlockLoop *loop)
\end{apient}

\apidesc{
  Find instrumentation points for the given loop that correspond to
  the given location: loop entry, loop exit, the start
  of a loop iteration and the end of a loop iteration.
  \code{BPatch\_locLoopEntry}  and \code{BPatch\_locLoopExit}  instrumentation
  points respectively execute once before the first iteration of a
  loop and after the last iteration.
  \code{BPatch\_locLoopStartIter}  and \code{BPatch\_locLoopEndIter}
  respectively execute at the beginning and end of each loop
  iteration.
}

\begin{apient}
  BPatch_basicBlock* findBlockByAddr(Dyninst::Address addr);
\end{apient}

\apidesc{
  Find the basic block within this flow graph that contains addr.
  Returns NULL on failure. This method is inefficient but
  guaranteed to succeed if addr is present in any block in this CFG.
  [NOTE: Dyninst is not always able to generate a correct flow graph
    in the presence of indirect jumps. If a function
    has a case statement or indirect jump instructions, the targets
    of the jumps are found by searching instruction
    patterns (peep-hole). The instruction patterns generated are
    compiler specific and the control flow graph analyses
    include only the ones we have seen. During the control flow graph
    generation, if a pattern that is not handled is used
    for case statement or multi-jump instructions in the function
    address space, the generated control flow graph may not
  be complete.]
}
